\section{Exact queries and independent unfoldings}
\label{sec:queries}

Fix a graph $G$ supplied by Proposition~\ref{prop:prefix}, and retain
$\cE=E(G)$ and the family $\cT$ of triangle edge sets.  We first express
edge survival by a finite recursive test.  We then give each occurrence
in that recursion fresh randomness.  The resulting auxiliary model has
exact product identities, which we will estimate before comparing it
with the finite process.

\subsection{Priorities and the finite test}

Independent priorities give the standard random greedy packing
representation; see Spencer~\cite[Section~1]{Spencer1995},
Grable~\cite[Section~1]{Grable1997}, and the triangle-removal accounts
\cite[Section~1.2]{BohmanFriezeLubetzky2015}
\cite[Section~14.1]{joosKuehn2025}.

Assign independent uniform priorities $U_T\in[0,1]$ to the triangles
$T\in\cT$.  Scan the triangles in increasing priority order, accepting a
triangle precisely when all three of its edges remain, and deleting
those edges on acceptance.  We work on the probability-one event that
all priorities are distinct.  An edge is \emph{untouched before $t$} if
no accepted triangle with priority strictly below $t$ contains it.

\begin{lemma}[Priority representation]\label{lem:priority}
For any finite simple graph $G$, the graph sequence indexed by the
accepted triangles in this scan has the law of uniform sequential
triangle removal started from $G$.
\end{lemma}

\begin{proof}
The priority order is a uniform random permutation of $\cT$.
Conditional on any scanned prefix, the unscanned triangles have uniform
relative order.  Every currently present triangle is unscanned: an
accepted triangle has lost its edges, and a rejected triangle already
had a deleted edge and can never become present again.  Until the next
acceptance the graph stays fixed.  Thus the next accepted triangle, if
there is one, is the first currently present triangle in the remaining
permutation, and is uniform among those triangles.  This also identifies
the absorbing state when none remains.
\end{proof}

In applying the lemma after the good prefix, the priorities are fresh
random variables on the triangles of the fixed graph $G$.  The original
process has the same continuation law because its next-step rule
depends only on its current graph.

A \emph{call} consists of a focus (one or two edges), a threshold $t$,
and, for a two-edge focus, a parent triangle.  The finite test has two
forms:
\begin{itemize}
\item A root call has focus $\{e\}$, threshold $t\in[0,1]$, and no
parent.  Its candidates are all triangles containing $e$.
\item A non-root call has focus $\{f,g\}\subset T$, parent triangle $T$,
and threshold $t=U_T$.  Its candidates are all triangles meeting the
focus, except $T$.
\end{itemize}
Distinct triangles meet in at most one graph edge.  Consequently every
candidate $S$ contains exactly one focus edge, denoted $e(S)$.  Its child
call has focus $S\setminus\{e(S)\}$, parent $S$, and threshold $U_S$.

At a call the test performs the following steps.
\begin{enumerate}
\item Expose the priorities of \emph{all} candidates, in a fixed
deterministic order of triangle types.  Each candidate occurs once in
this list.
\item Sort the candidates by priority and consider those with $U_S<t$
in that order.  For a pair focus this is one joint ordering of the
candidates at both focus edges.
\item Evaluate each considered candidate's child.  At the first child
returning true, stop and return false.  If no child returns true, return
true.
\end{enumerate}
Exposing a candidate does not require visiting its child: priorities
above the threshold and candidates skipped after an earlier success
are also exposed.  All calls use the same finite family $(U_T)_{T\in\cT}$.
The test has no side effects; it does not delete edges or change any
priority.

\begin{lemma}[Exact paired test]\label{lem:finite-test}
For any finite simple graph with distinct triangle priorities, the
finite test terminates.  A root call at $(e,t)$ returns true exactly when
$e$ is untouched before $t$ in the priority scan.  A call with focus
$\{f,g\}\subset T$ and threshold $U_T$ returns true exactly when both
$f$ and $g$ are untouched before $U_T$.
\end{lemma}

\begin{proof}
Along a recursive path the thresholds strictly decrease through the
priorities of triangle types.  No type can occur twice on that path,
so its length is at most $|\cT|$.  The branching is finite, proving
termination.

We prove correctness by induction on the number of triangle priorities
strictly below the call's threshold.  Every child has a smaller
threshold, so the induction hypothesis applies to it.  Write $\mathcal F$
for the focus of the present call.  We claim that some edge of
$\mathcal F$ is deleted before $t$ if and only if some candidate $S$
with $U_S<t$ has both edges of $S\setminus\{e(S)\}$ untouched before
$U_S$.

For the forward implication, take an accepted triangle that deletes a
focus edge before $t$.  It is a candidate and all its edges are
untouched just before its acceptance.  The omitted parent, when
present, is at the threshold itself and cannot be this triangle.
Conversely, suppose such a candidate exists and every focus edge is
untouched before $t$.  Its focus edge is then present at $U_S$, as are
its two other edges.  The scan accepts it, contradicting the assumed
survival of the focus.  This proves the claim.  By induction the
candidate condition is precisely that its child returns true, and the
test returns false exactly when at least one such child exists.
\end{proof}

The parent omission is local to its call and is justified by the strict
threshold.  It does not erase that triangle from descendant lists.
Also, the two finite survival events in a pair call need not be
independent.  We retain their joint chronological test, including its
stopping rule, throughout the comparison argument.

\subsection{Fresh randomness at each occurrence}

Backward exploration and comparison with independent branching occur in
Spencer's analysis \cite[Sections~2--4]{Spencer1995}. We retain the paired
query rule above and all candidate-list exposures for the finite comparison.

For a fixed root type, first construct the deterministic tree of all
possible calls.  Each node has one child for each candidate on its
list, with the focus and parent specified above.  The tree is countable
and finitely branching.  A triangle type may label many different
child occurrences.

Assign an independent uniform variable to each child occurrence.  This
variable serves both as the candidate priority in its parent's list
and as the child's threshold if the child is visited.  Apply the same
joint scanning and stopping rule on this tree.  Different occurrences
of a triangle type now have independent priorities; different root
queries use disjoint trees with independent variables.  We call this
the \emph{independent unfolding}.

We also allow a singleton root of type $e$ with a specified triangle
$T\ni e$ omitted from its root list.  The omission affects that list
only.  Descendant occurrences of $T$ remain available whenever their
own parent rule allows them.  A two-edge root with focus
$\{f,g\}\subset T$, omitted parent $T$, and prescribed threshold $t$
is defined by the same independent construction.

The finiteness argument uses factorial decay along decreasing-priority
paths, as in the local construction of graph adsorption processes by
Penrose and Sudbury~\cite[Proposition~1 and its proof]{PenroseSudbury2005}.
We give the argument for the present occurrence trees.

\begin{lemma}[Finite evaluation]\label{lem:unfolding-finite}
For a fixed good graph $G$, independent unfoldings can be constructed
so that, almost surely, every possible call and all thresholds
$t\in[0,1]$ have finite evaluations.  Their outputs are jointly
measurable in the threshold and the occurrence priorities.
\end{lemma}

\begin{proof}
Put $d_{\max}=\max_{e\in\cE}d_e$ and $M=2d_{\max}$; every candidate
list has at most $M$ entries.  For large $n$, the good-graph estimates
give $M\le4D$.  From any fixed possible node there are at most $M^j$
descendant paths of length $j$.  On each path the $j$ occurrence
priorities are independent uniforms, and the probability that they
strictly decrease is $1/j!$.  Hence
\[
 \Prb(\text{a decreasing path of length $j$ starts at that node})
 \le \frac{M^j}{j!}\longrightarrow0.
\]
The events on the left decrease with $j$.  Almost surely their depths
are therefore bounded; finite branching makes the entire decreasing
dependency tree finite.  This argument uses threshold $1$, and thus
works simultaneously for every smaller threshold.  There are only
countably many nodes in the forests for all possible root types, so a
countable intersection gives the asserted simultaneous conclusion.

Truncate the recursion at depth $j$, using any fixed Boolean value at
the boundary.  Its output is a measurable function of the threshold
and finitely many priorities.  On the probability-one event just
proved, these outputs stabilize as $j\to\infty$, for all thresholds,
to the output of the full evaluation.  Defining an arbitrary value on
the null exceptional set proves joint measurability.
\end{proof}

We may consequently evaluate a child's subtree hypothetically even
when the chronological test never reaches that child.  These
hypothetical outputs characterize the root event: a call returns true
if and only if no child with priority below its threshold would return
true.  For fixed priorities this indicator is nonincreasing in the
root threshold, since each child's threshold is its own occurrence
priority.

\subsection{Exact probability identities}

Let $q_e(t)$ be the probability that a full singleton root at $e$
returns true at threshold $t$ in the independent unfolding.  Let
$q_{e,T}(t)$ be the same probability when $T\ni e$ is omitted only
from that singleton root list.  For $T=\{e,f,g\}\in\cT$, define
\begin{equation}\label{eq:integral}
 I_{e,T}(t)=\int_0^t q_{f,T}(v)q_{g,T}(v)\,dv,
 \qquad 0\le t\le1.
\end{equation}
The integral is well defined by Lemma~\ref{lem:unfolding-finite}.

\begin{proposition}[Product identities]\label{prop:products}
For every triangle $T=\{e,f,g\}\in\cT$ and $t\in[0,1]$, the
independent unfolding satisfies
\begin{equation}\label{eq:pair}
 \Prb(\text{the pair $\{f,g\}\subset T$ returns true at $t$})
   =q_{f,T}(t)q_{g,T}(t),
\end{equation}
\begin{equation}\label{eq:products}
 \begin{aligned}
 q_e(t)&=\prod_{S\ni e}\bigl(1-I_{e,S}(t)\bigr),\\
 q_{e,T}(t)&=\prod_{\substack{S\ni e\\S\ne T}}
                    \bigl(1-I_{e,S}(t)\bigr).
 \end{aligned}
\end{equation}
In the pair probability, $T$ is omitted from the pair's candidate list.
All products are over $\cT$, and an empty product equals $1$.
The functions $q_e$, $q_{e,T}$, and $I_{e,T}$ belong to $C^1([0,1])$.
Moreover, $q_e(0)=q_{e,T}(0)=1$ and $I_{e,T}(0)=0$.
\end{proposition}

\begin{proof}
After omitting $T$, the candidate sets at $f$ and $g$ are disjoint:
a triangle containing both edges would equal $T$.  Their occurrence
variables and all descendant variables are independent.  The pair
returns true precisely when every candidate in both sets fails to
give a successful child below $t$.  The two separate events have the
laws of the singleton roots with $T$ omitted, proving
\eqref{eq:pair}.  The ordering of the joint list changes which children
are visited, but not this event that all relevant children fail.

Consider a candidate $S=\{e,f,g\}$ at a singleton root $e$.
Conditional on its occurrence priority being $v$, its child is a pair
call at threshold $v$, with $S$ omitted and independent descendant
priorities.  Equation~\eqref{eq:pair} gives its success probability as
$q_{f,S}(v)q_{g,S}(v)$.  Thus $I_{e,S}(t)$ is exactly the probability
that this candidate has priority below $t$ and a successful child.
Different root candidates use disjoint sets of occurrence variables,
including their descendants.  Taking the probability that none
succeeds proves both products in \eqref{eq:products}.

Initially the probabilities are measurable and bounded between $0$
and $1$.  Equation~\eqref{eq:integral} therefore makes each $I_{e,T}$
Lipschitz.  The finite products imply continuity of all $q_e$ and
$q_{e,T}$.  Substitution back into \eqref{eq:integral} now gives a
continuous integrand, so the integrals, and then the products, are
continuously differentiable.  Derivatives at the endpoints have their
one-sided interpretation.  The initial values follow directly from
the formulas.
\end{proof}

These identities describe the auxiliary independent model.  They
assert no factorization of survival probabilities in the finite
priority system.  We next estimate their solution uniformly over all
edge and omitted-parent types; the finite evaluation above has already
constructed this solution, so no recursive fixed-point assumption is
needed.
